% EM structure and teaching revision 20260928
\subsection{磁力做功与电源供能如何相容}
单个电荷的磁力垂直于它的总速度，不做功；载流导体却可以受到安培力并获得机械功。这两句不矛盾，因为导体平动、载流子的相对漂移以及维持电流的电源共同参与能量交换。
长为 $\ell$ 的导体棒在匀强磁场中垂直于棒和磁场运动，取电流方向使磁力推动棒，大小为 $F=BI\ell$。棒的运动产生反电动势 $\mathcal E_b=B\ell v$，因而
\[
 P_{\mathrm{mech}}=Fv=\mathcal E_b I.
\]
若电源电动势为 $\mathcal E_0$、总电阻为 $R$，稳定电流满足
$\mathcal E_0=IR+\mathcal E_b$，电源功率分为焦耳热与机械输出：
$\mathcal E_0I=I^2R+Fv$。
外力反过来拖动棒时，机械输入又可变为电能，进入电磁感应问题。判断功率方向应统一电流和速度的参考方向，不应只因出现磁力就断言“体系没有机械能变化”。

